\documentclass[14pt,usepdftitle=false,aspectratio=169]{beamer}
\usepackage{preambule}
\setbeamercolor{structure}{fg=black}
\usepackage{polynomes,al}
\hypersetup{pdftitle={Théorie spectrale -- Spectre dans une algèbre quelconque}}
\title{Théorie spectrale\\\emph{Spectre dans une algèbre quelconque}}
\author{}
\date{}
\begin{document}
\begin{frame}
    \titlepage
\end{frame}
\slideq{CNS pour que $P\l a\r\in A^\times$ sur $\bbK$ algébriquement clos}{1/8}
\slider{$P$ ne s'annule pas sur $\sigma\l a\r$}{1/8}
\slideq{Spectre de $a\in A$ avec $A$ une $\bbK$-algèbre}{2/8}
\slider{$\sigma_A\l a\r=\set{\lambda\in\bbK,a-\lambda\in A^\times}$}{2/8}
\slideq{Propriétés topologiques de $A^\times$ pour $A$ une algèbre de Banach}{3/8}
\slider{$A^\times$ est un ouvert et est un groupe topologique}{3/8}
\slideq{Résolvante de $a$}{4/8}
\slider{$\appl{R_a}{\bbC\setminus\sigma\l a\r}{A^\times}{\lambda}{\l a-\lambda\r^{-1}}$\linebreak C'est une application holomorphe et $R_a'=R_a^2$}{4/8}
\slideq{$F\l a\r$ pour $F\in\fr KX$}{5/8}
\slider{Si $F=\frac PQ$ avec $P,Q\in\pol KX$ et $Q$ n'a pas de pôles sur $\sigma\l a\r$\linebreak$F\l a\r=P\l a\r Q\l a\r^{-1}$}{5/8}
\slideq{Lien entre $\sigma\l ab\r$ et $\sigma\l ba\r$}{6/8}
\slider{$\sigma\l ab\r\setminus\set0=\sigma\l ba\r\setminus\set0$}{6/8}
\slideq{$\sigma\l F\l a\r\r$ pour $F\in\fr KX$}{7/8}
\slider{Si $F$ est non constante et n'a pas de pôles sur $\sigma\l a\r$\linebreak$\sigma\l F\l a\r\r=F\l\sigma\l a\r\r$}{7/8}
\slideq{Propriétés topologiques de $\sigma\l a\r$ pour $a\in A$ une algèbre de Banach}{8/8}
\slider{C'est un compact non vide contenu dans la boule de rayon $\left\|a\right\|$}{8/8}
\end{document}